\documentclass[11pt]{article}
\usepackage[margin=1in]{geometry}
\usepackage{amsmath,amssymb,amsthm,enumitem,booktabs,array}
\newtheorem{theorem}{Theorem}
\newtheorem*{claim}{Claim}
\title{Proof check request: Theorem 7.4\\[2pt]\large ``Stress identifies the leverage distribution''}
\author{Pratik Niroula, Prashant Niroula and Sujan Khatiwada}
\date{10 November 2026}
\begin{document}
\maketitle

\noindent\textbf{What we ask.} Please check two things and fill in the table in Section~5.
\begin{itemize}[itemsep=1pt]
\item Does the Lean statement in Section~2 say what the prose theorem in Section~1 says?
\item Is the one part proved only on paper, the small-probe statement, correct (Section~3)?
\end{itemize}
The theorem is new. It is Theorem~7.4 of our working paper, ``A Sufficient Statistic for Algorithmic Feedback''. An AI system (Claude) drafted the statement, the proofs and the Lean code. The Lean proofs compile against Mathlib with the three standard axioms only, so the question for a human reader is no longer whether the main argument has a gap. It is whether the formal statement is the right one.

\section{Setting and statement}

\paragraph{Market.} Prices follow $x_t=ax_{t-1}+\kappa(F_t+v_t)$ for $t\ge0$, with $x_{-1}=0$ and $\kappa\ne0$. Here $v_t$ is a probe sequence chosen by the experimenter, and the response $(x_t)$ is observed exactly.

\paragraph{Traders.} There are finitely many traders. Trader $i$ has a memory rate $r_i\ge0$, a capacity $c_i$ and a leverage $\ell_i\ne0$. Its memory state is $m_{i,t}=\sum_{j<t}r_i^{\,t-1-j}x_j$, and its flow is $c_i\psi(\ell_im_{i,t})$. Aggregate flow is $F_t=\sum_ic_i\psi(\ell_im_{i,t})$, for an odd shape $\psi$ common to all traders.

\paragraph{Objects.} The signed capacity measure is $\tilde\mu=\sum_ic_i\,\mathrm{sgn}(\ell_i)\,\delta_{(r_i,|\ell_i|)}$. A shape has \emph{independent dilations} if, for distinct $s_1,\ldots,s_n>0$, $\sum_kw_k\psi(s_ku)=0$ for every real $u$ implies every $w_k=0$.

\begin{theorem}
Let $\psi$ be odd with independent dilations. Two populations give identical responses to every probe sequence if and only if $\tilde\mu_P=\tilde\mu_{P'}$.
\begin{enumerate}[label=(\roman*)]
\item The hard limit $\mathrm{clip}(u)=\max(-1,\min(1,u))$ has independent dilations.
\item $\tanh$ has independent dilations. Moreover, identical responses to all probes in an arbitrarily small neighbourhood of zero already imply $\tilde\mu_P=\tilde\mu_{P'}$.
\item A cubic shape $u-\gamma u^3$ does not have independent dilations.
\end{enumerate}
\end{theorem}

\section{What Lean checks}

The following is the type of the formal theorem for $\tanh$, as printed by Lean (\texttt{p1\_thm\_stress\_identification\_tanh}). The general theorem and the hard-limit theorem have the same form.
{\small
\begin{verbatim}
forall {K K' : Nat} (a kappa : Real), kappa != 0 ->
  forall (c l r : Fin K -> Real) (c' l' r' : Fin K' -> Real),
    (forall i, 0 <= r i) -> (forall j, 0 <= r' j) ->
    (forall i, l i != 0) -> (forall j, l' j != 0) ->
    ((forall (v x : Nat -> Real),
        ProbePricePath a kappa (satAggFlow tanh c l r) v x ->
        ProbePricePath a kappa (satAggFlow tanh c' l' r') v x)
     <->
     forall (g : Real -> Real -> Real),
       Sum_i c i * realSgn (l i) * g (r i) |l i|
         = Sum_j c' j * realSgn (l' j) * g (r' j) |l' j|)
\end{verbatim}}
\noindent Here \texttt{ProbePricePath a kappa F v x} says $x_0=\kappa(F_0(x)+v_0)$ and $x_{t+1}=ax_t+\kappa(F_{t+1}(x)+v_{t+1})$ for all $t$; \texttt{satAggFlow psi c l r x t} is $\sum_ic_i\psi(\ell_im_{i,t})$ with $m_{i,0}=0$ and $m_{i,t+1}=r_im_{i,t}+x_t$.

\paragraph{Points where the formal statement differs from the prose.}
\begin{itemize}[itemsep=1pt]
\item ``Identical responses'' is formalized as: every (probe, path) pair of the first market is a pair of the second. Uniqueness of the path for a given probe is not used and not proved.
\item Equality of the measures is formalized as equality of the integrals of every function $g$ of (memory, leverage).
\item Capacities $c_i$ are arbitrary reals in Lean. The prose has $c_i>0$.
\item Independent dilations quantifies over all real $u$, so the Lean theorems use probes of every size.
\end{itemize}

\paragraph{Proof outline (all machine-checked).} Probes steer the price along any path, so identical responses give equal flow maps. On the path $x_0=\theta u$, $x_1=u$, the date-2 flow is $\sum_kd_k\psi(s_ku)$ with $s_k=\lambda_k(r_k\theta+1)$, and a direction $\theta>0$ exists for which distinct atoms have distinct scales. Independent dilations then give $d_k=0$. For the hard limit, the largest scale is isolated at $u=1/s_1$ and $u=1/s_2$. For $\tanh$, $1-\tanh y=2e^{-2y}/(1+e^{-2y})$, so the smallest scale has the slowest tail and is isolated by multiplying by $e^{2s_1u}$ and letting $u\to\infty$.

\section{The part proved only on paper: small probes}

\begin{claim}
If $g(u)=\sum_kw_k\tanh(s_ku)=0$ for all $u$ in a neighbourhood of $0$, with distinct $s_k>0$, then every $w_k=0$.
\end{claim}
\noindent\emph{Proof A.} $g$ is real-analytic on $\mathbb R$, so it vanishes on all of $\mathbb R$, and the machine-checked result applies.

\noindent\emph{Proof B.} $\tanh u=\sum_{n\ge1}\tau_{2n-1}u^{2n-1}$ with $\tau_{2n-1}=2^{2n}(2^{2n}-1)B_{2n}/(2n)!\ne0$. So $\sum_k(w_ks_k)(s_k^2)^{n-1}=0$ for $n=1,\ldots,K$, a Vandermonde system in the distinct numbers $s_k^2$.

\noindent For the theorem, small target paths need only small probes, because $|F_t|\le\sum_i|c_i||\ell_i||m_{i,t}|$.

\section{Numerical check (not a proof)}

The script \texttt{research\_extension/nonlinear\_scoping/thm74\_numeric\_check.py} (also packaged as the Colab notebook \texttt{Thm74\_Check\_Colab.ipynb}) runs five checks for three seeds:
\begin{itemize}[itemsep=1pt]
\item \textbf{T1.} Same $\tilde\mu$, after random splits and offsetting pairs, gives identical responses. Largest gap: $1.4\times10^{-14}$ over 40 populations and 200 random probes, under both shapes.
\item \textbf{T2.} Same kernel, different $\tilde\mu$: under clip, identical below pulse size $5$, then separated. Under $\tanh$, the first-tap gap scales with fitted exponent $2.98$, against the predicted $3$.
\item \textbf{T3.} Constructive recovery of $\tilde\mu$ under clip from two directions recovers all atoms in 25 of 25 random populations. Atoms were first matched by exact rounded keys, and 3 cases failed by one unit in the third decimal (the grid step). With a tolerance equal to the grid step, all 25 match.
\item \textbf{T4.} Cubic Taylor coefficient of the $\tanh$ flow-map difference: fitted $9.2264061\times10^{-4}$, predicted $9.2264063\times10^{-4}$.
\item \textbf{T5.} The counterexample to the old wording of (ii): the gap is $0$ inside all bands and $0.375$ between the two band edges.
\end{itemize}

\section{Questions and response table}

\begin{enumerate}[itemsep=2pt]
\item Does the Lean statement in Section~2 match the prose theorem, given the four listed differences? Is any of them a real weakening?
\item Is ``every pair of the first market is a pair of the second'' an acceptable reading of ``identical responses to every probe''?
\item Is the definition of independent dilations the right hypothesis? Is there a natural weaker one?
\item Is the small-probe claim correct, by Proof A or Proof B?
\item Is $\tilde\mu$ exactly the equivalence class for odd shapes with independent dilations?
\item Do you know this result, or the dilation-independence facts for clip and $\tanh$, from the literature on ridge functions or neural-network identifiability? A reference would help us position it.
\end{enumerate}

\vspace{4pt}
\begin{tabular}{>{\raggedright\arraybackslash}p{0.12\textwidth}>{\raggedright\arraybackslash}p{0.2\textwidth}>{\raggedright\arraybackslash}p{0.56\textwidth}}
\toprule
Question & Correct / gap / error & Comment\\
\midrule
1 & & \\[14pt]
2 & & \\[14pt]
3 & & \\[14pt]
4 & & \\[14pt]
5 & & \\[14pt]
6 & & \\[14pt]
\bottomrule
\end{tabular}

\vspace{6pt}
\noindent Reader: \rule{4.5cm}{0.4pt} \quad Date: \rule{3cm}{0.4pt}\quad We will acknowledge the reader by name in the paper only with their permission.
\end{document}
